\reviewsection{Media Observation to Confirm Before Submission}

\begin{center}
\colorbox{SoftGold}{%
\parbox{0.92\textwidth}{%
\small
\textbf{Required review insertion:} After completing the selected Module 3 media, add one verified scene, statement, or interaction here. Name what the communicator did, what the partner or environment did, and whether the response increased or restricted agency. The value of the evidence does not depend on whether the source was accessed audiovisually, through captions, or through an accessible transcript; it depends on connecting an accurate, concrete moment to the module's essential question.
}%
}
\end{center}

\reviewsection{What I Will Carry Into My Classroom}

In my future classroom, I will treat AAC as a multimodal system rather than one device. Students need vocabulary for choosing, refusing, commenting, asking, joking, participating academically, expressing sensory needs, and repairing misunderstandings. They need primary and backup options that travel through arrival, instruction, transitions, group work, lunch, and dismissal. Adults and peers need to model, wait, notice, confirm, and act without requiring performance or taking over the student's voice. The team must also notice when a gesture, facial expression, object, written message, or pause already communicates effectively rather than demanding device repetition to make the message look more conventional \citep{mirenda2015aids}.

Meaningful communication therefore looks like more than access to words. It looks like a student being recognized as the author of a message that matters. It looks like partners who can tolerate being corrected and plans that change when the communicator shows that a support is not working. It looks like family knowledge traveling into school and student-authored information traveling home. Most of all, it looks like a learner having an accessible way to challenge an interpretation, enter a relationship, and change what happens next.
